\section{Finite realified normal form}
\label{sec:normal-form}

\subsection{Spaces and pairing}

Let \(I\) and \(J\) be finite index sets. Regard
\(\C^I,\C^J\) as real vector spaces and use
\[
 [x,y]_{\mathbb R}
 =
 \Real\langle x,y\rangle
 =
 \Real\sum_i\overline{x_i}y_i.
\]
The adjoint of a real-linear map \(C\) is characterized by
\[
 [C\theta,z]_{\mathbb R}
 =
 [\theta,C^*z]_{\mathbb R}.
\]
Complex-linear provenance filters are the main case, while the
real-linear convention also covers a shared filter constrained by
conjugation.

For a real-linear map, \(C_{ij}\) below denotes the real-linear
coordinate block from the \(j\)-th copy of \(\C\) to the \(i\)-th copy,
not necessarily multiplication by one complex scalar. Thus
\(C_{ij}\ne0\) means that this block map is nonzero. In the
complex-linear pair-dependent program, it is the usual scalar matrix
entry.

\begin{definition}[Affine completion value]
\label{def:Gamma}
For \(b\in\C^I\) and a real-linear map
\(C:\C^J\to\C^I\), define
\[
 \Gamma(b,C)
 =
 \min_{\theta\in\C^J}\|b+C\theta\|_1.
\]
\end{definition}

The minimum exists. Indeed, \(b+\Ran C\) is a closed affine subspace in
finite dimension, and a nearest point to the origin exists. The
parameter \(\theta\) need not be unique.

\subsection{Global dual}

Finite-dimensional distance-to-a-subspace duality
\citep[Section 15]{Rockafellar1970} gives
\begin{equation}
 \Gamma(b,C)
 =
 \max_{\substack{\|z\|_\infty\le1\\C^*z=0}}
 [z,b]_{\mathbb R}.
 \label{eq:global-dual}
\end{equation}
Both extrema are attained. The coordinatewise unit polydisc, rather
than a real cube, is the dual ball of complex \(\ell^1\).

\subsection{Exact support}

All graph statements use
\[
\supp C=\{(i,j)\in I\times J:C_{ij}\ne0\}
\]
after exact assembly. If entries are validated intervals, an edge may
be deleted only when its interval is the singleton \(\{0\}\). A small
floating magnitude is not an exact zero.

\begin{remark}[Canonical coordinates]
Changing bases in parameter space can alter the support graph while
preserving \(\Ran C\). The graph factorization below concerns the
declared literal filter coordinates. It is maximal among decompositions
read directly from that coordinate support, not among all possible
algebraic decompositions of the quotient.
\end{remark}
